\section{生体分子と高分子の構造}\label{theory:organic:biomolecules}
アミノ酸はアミノ基とカルボキシ基を持つ化合物である。タンパク質を作るアミノ酸では側鎖の種類が多様で、OHを持つものもあるが、遊離水酸基は全アミノ酸の必須条件ではない\cite{analysis49_amino_ref}。
アミノ酸残基を結ぶペプチド結合は主鎖のアミド結合である。鎖の共有結合順序と、二次・三次構造の配置を分ける。変性では主として高次構造が変わり機能が失われ得るが、必ず主鎖が加水分解されるわけでも全て不可逆なわけでもない。
等電点 \(pI\) では正味電荷が零となる。pH勾配を使う等電点電気泳動では、pIの異なるタンパク質を異なる位置へ集められる。

酵素の多数はタンパク質だが、触媒RNAであるリボザイムも知られる。1989年のNobel化学賞はRNAの触媒作用の発見を対象としている\cite{analysis49_rna_ref}。平成27・30年度の「すべてタンパク質」という原典肢を、公式の別肢選択だけで真と確定しない。

高分子の分類では天然・合成、熱可塑性・熱硬化性、単独重合・共重合、混合を別々に考える。通常のPEは熱可塑性、ナイロンは合成ポリアミドの代表例である\cite{organic_physical,organic_step}。
セルロースは多数のOHを持っても、鎖同士の相互作用と凝集・結晶性等を含む溶解の自由エネルギーにより通常水へ難溶である\cite{analysis49_cellulose_ref}。局所官能基の親水性だけから巨大な鎖全体の易溶性を推論しない。水素結合は重要だが、それだけで全溶解挙動を説明できるとはしない。
